\documentclass[conference]{IEEEtran}

\usepackage{amsmath}
\usepackage{graphicx}
\usepackage{cite}

\title{In-body Antigen Detection via Antibodies on a Transistor}
\author{Duran Kaan Altin}

\begin{document}
\maketitle

\begin{abstract}
This report studies a COMSOL-based graphene field-effect transistor biosensor model for HER2 detection. The workflow couples HER2 transport, antibody surface binding, GFET current response, sensitivity analysis, and limit of detection estimation.
\end{abstract}

\section{Introduction}
Describe the HER2 detection problem and the motivation for label-free GFET biosensing.

\section{Background}
Summarize HER2 biomarkers, antigen-antibody binding, field-effect biosensors, and GFET sensing.

\section{Methodology}
Describe the kinetic model, transport model, surface binding boundary, and electrical response mapping.

\section{Results}
Present concentration profiles, bound HER2 density, current response curves, sensitivity, and LOD.

\section{Discussion}
Interpret whether response is transport-limited, electronics-limited, or both under the selected assumptions.

\section{Limitations and Future Work}
Discuss wet-lab validation, Debye screening, antibody immobilization, COMSOL approximations, and implantability.

\section{Conclusion}
Summarize the final modeling contribution and main engineering insight.

\bibliographystyle{IEEEtran}
\bibliography{references}

\end{document}
